\documentclass[11pt,a4paper]{article}
\usepackage[T1]{fontenc}
\usepackage{lmodern}
\usepackage[margin=22mm,headheight=14pt]{geometry}
\usepackage{amsmath,amssymb,graphicx,booktabs,tabularx,array}
\usepackage{xcolor,listings,enumitem,microtype,fancyhdr}
\usepackage{tikz}
\usetikzlibrary{arrows.meta,positioning}
\usepackage[font=small,labelfont=bf]{caption}
\captionsetup{hypcap=false}
\usepackage[hidelinks]{hyperref}
\hypersetup{pdftitle={MIT VNAV Individual Deliverables},
pdfauthor={yawangjiang}}
\definecolor{navy}{HTML}{19394F}
\definecolor{pale}{HTML}{F0F4F7}
\setlength{\parindent}{0pt}
\setlength{\parskip}{5pt}
\setlength{\emergencystretch}{2em}
\setlist[itemize]{leftmargin=1.5em,itemsep=3pt,topsep=3pt}
\setlist[enumerate]{leftmargin=1.7em,itemsep=3pt,topsep=3pt}
\pagestyle{fancy}
\fancyhf{}
\fancyhead[L]{\small\color{navy}MIT VNAV - Individual}
\fancyhead[R]{\small\color{navy}yawangjiang}
\fancyfoot[C]{\small\thepage}
\renewcommand{\headrulewidth}{0.3pt}
\lstset{basicstyle=\ttfamily\small,columns=fullflexible,keepspaces=true,
breaklines=true,breakatwhitespace=false,showstringspaces=false,
backgroundcolor=\color{pale},frame=single,rulecolor=\color{pale},
xleftmargin=5pt,xrightmargin=5pt,framesep=5pt,aboveskip=7pt,belowskip=7pt}
\newcommand{\code}[1]{\texttt{#1}}
\newcommand{\R}{\mathbb{R}}
\newcommand{\ex}{\mathbf e_x}
\newcommand{\ev}{\mathbf e_v}
\newcommand{\eR}{\mathbf e_R}
\newcommand{\eO}{\mathbf e_\Omega}
\newcommand{\eThree}{\mathbf e_3}
\newcommand{\Fdes}{\mathbf F_{\mathrm{des}}}
\newcommand{\Fdrag}{\mathbf F_{\mathrm{drag}}}
\newcommand{\repo}{https://github.com/LittleHorseyaya/vnav}

\hypersetup{pdftitle={MIT VNAV LAB3 - Individual Deliverable 2},pdfauthor={yawangjiang}}
\begin{document}
\begin{center}
{\LARGE\bfseries\color{navy}LAB3: Modelling and Control of UAVs}\par
\vspace{6pt}
{\large Individual Deliverable 2}\par
\vspace{6pt}
Submission account: \code{yawangjiang}\par
October 31, 2026
\end{center}

\section{A. Structure of Quadrotors}
Let $c_f>0$ and $c_d>0$ be the lift and yaw-reaction coefficients, $d>0$ the
arm length, and $a=d c_f/\sqrt2$. Use an upward-$z$ body frame and the motor
numbering for which the functional vehicle has the following wrench map:
\begin{equation}
\begin{bmatrix}f\\\tau_x\\\tau_y\\\tau_z\end{bmatrix}
=F_a\mathbf u,\qquad
F_a=\begin{bmatrix}
c_f&c_f&c_f&c_f\\
a&a&-a&-a\\
-a&a&a&-a\\
c_d&-c_d&c_d&-c_d
\end{bmatrix}.
\label{eq:fa}
\end{equation}
Here $u_i=\omega_i^2\geq0$ describes the positive lift-producing operating
direction of each mounted propeller. The alternating yaw signs already encode
the different positive physical spinning directions.

For vehicle (a),
\begin{equation}
\det(F_a)=-16a^2c_fc_d\ne0,\qquad\operatorname{rank}(F_a)=4.
\end{equation}
Thus small changes of collective thrust and the three torques can be independently
allocated around a feasible hover. Full rank alone does not guarantee feasibility
for arbitrarily large commands under motor bounds.

In vehicle (b), propellers 3 and 4 are changed and their spin directions are
reversed. They continue to produce upward lift, but their yaw-reaction signs
reverse. With the same numbering and frame,
\begin{equation}
F_b=\begin{bmatrix}
c_f&c_f&c_f&c_f\\
a&a&-a&-a\\
-a&a&a&-a\\
c_d&-c_d&-c_d&c_d
\end{bmatrix}.
\end{equation}
Its fourth row is $-c_d/a$ times its third row. The first three rows remain
independent, so
\begin{equation}
\boxed{\operatorname{rank}(F_b)=3,\qquad
\tau_z=-\frac{c_d}{a}\tau_y=-\frac{\sqrt2c_d}{dc_f}\tau_y.}
\end{equation}
Therefore yaw torque is coupled to pitch torque for this convention. A different
documented motor order can change the row signs or the named horizontal torque,
but not the rank deficiency.

For example, consider a level hover with zero angular velocity and a reference
that holds position while requesting a yaw angular acceleration. Nonzero yaw
torque necessarily produces the linked pitch torque, which tilts the thrust
axis and disturbs horizontal position. Consequently vehicle (b) cannot generally
track arbitrary independent references for $[x,y,z]$ and yaw. Some specially
compatible motions can still be achieved.

\clearpage
\section{B. Compensation of the Specified Drag}
The empirical body-frame force in the exercise is quadratic in each velocity
component and acts against that component. Let
\begin{equation}
D=\operatorname{diag}(0.1,0.1,0.2),\qquad
\mathbf v^b=R^\top\mathbf v,\qquad
\mathbf F^b_{\mathrm{drag}}=-D\bigl(\mathbf v^b\odot|\mathbf v^b|\bigr),
\end{equation}
where $R$ maps body vectors to world vectors, $\odot$ is the componentwise
product, and the absolute value is componentwise. A linear expression
$-D\mathbf v^b$ would describe a different force model.

The world-frame disturbance is
\begin{equation}
\mathbf F_{\mathrm{drag}}=R\mathbf F^b_{\mathrm{drag}}.
\end{equation}
With upward $z$, the translational dynamics are
\begin{equation}
m\dot{\mathbf v}=-mg\mathbf e_3+fR\mathbf e_3+\mathbf F_{\mathrm{drag}}.
\end{equation}
Define $\mathbf e_x=\mathbf x-\mathbf x_d$ and
$\mathbf e_v=\mathbf v-\mathbf v_d$. If $K_x,K_v$ are acceleration-feedback
gains, the nominal commanded force is
\begin{equation}
\mathbf F_0=m\left(\mathbf a_d+g\mathbf e_3-K_x\mathbf e_x-K_v\mathbf e_v\right).
\end{equation}
The model can be explicitly compensated by subtracting the estimated disturbance:
\begin{align}
\mathbf F_{\mathrm{comp}}
&=\mathbf F_0-\widehat{\mathbf F}_{\mathrm{drag}}\\
&=\mathbf F_0+RD\bigl(\mathbf v^b\odot|\mathbf v^b|\bigr)
\quad\text{under the stated model.}
\end{align}
Use this compensated force consistently in both the desired thrust direction
and scalar thrust projection:
\begin{equation}
\mathbf b_{3d}=\frac{\mathbf F_{\mathrm{comp}}}{\|\mathbf F_{\mathrm{comp}}\|},
\qquad f=\mathbf F_{\mathrm{comp}}^\top R\mathbf e_3.
\end{equation}
The remaining desired body axes are constructed from $\mathbf b_{3d}$ and the
desired yaw, as in the geometric controller.

As a sign check, take $R=I$ and $\mathbf v=[2,-1,0]^\top$\,m/s. Then
\begin{equation}
\mathbf F_{\mathrm{drag}}=[-0.4,0.1,0]^\top\ \mathrm N,
\qquad-\mathbf F_{\mathrm{drag}}=[0.4,-0.1,0]^\top\ \mathrm N.
\end{equation}
The commanded correction opposes the disturbance. In the ideal force-tracking
limit with an exact drag model, the extra dynamics term is cancelled, reducing
the disturbance that feedback must reject. Attitude tracking, model errors,
sampling delay, and actuator limits can still produce nonzero tracking error.
This is a modelling proposal, not a claim of an experimentally measured
before/after improvement.

\subsection*{Sources}
MIT VNAV LAB3 exercise statement:
\url{https://vnav.mit.edu/labs_2023/lab3/exercises.html}.\par
T. Lee, M. Leok, and N. H. McClamroch, ``Geometric tracking control of a quadrotor
UAV on SE(3),'' \emph{IEEE CDC}, 2010.
\end{document}
